\documentclass[a4paper,11pt]{article}
\usepackage{exam-style-341}
\examinehead{341 农业综合知识三 · 程序设计分章练习}{第04章参考答案}
\begin{document}

\unitanswerscover{341 农业综合知识三 · 程序设计分章练习}{第04章参考答案}{选择结构：逐题答案、解析与综合题完整参考实现；判断题均给出改正后的表述。}

\answerpart{一、单项选择题（10 题，每题 2 分，共 20 分）}

\q A。\begin{analysis}关系运算符 \code{<} 判断左操作数是否小于右操作数，3 $<$ 4 成立，关系表达式的值为 1（真）。若写成 \code{a > b} 则值为 0。\end{analysis}

\q B。\begin{analysis}\code{if (x = 5)} 的条件部分是赋值表达式而不是关系表达式：它先把 5 赋给 \code{x}，整个赋值表达式的值为 5。C 语言中非 0 为真，所以条件成立。这是把 \code{=} 误写成 \code{==} 的典型错误。\end{analysis}

\q B。\begin{analysis}\code{a++ \&\& a++} 从左到右计算：左边 \code{a++} 先取 \code{a} 的值 0 作为左操作数，\code{a} 自增为 1；由于逻辑与的左操作数为假，发生短路，右边的 \code{a++} 不再执行，\code{a} 仍为 1，整个表达式的值为 0，\code{if} 条件不成立，不输出 T。最后 \code{printf("\%d", a)} 输出 1。\end{analysis}

\q B。\begin{analysis}\code{k = 1} 与 \code{case 1} 匹配，执行 \code{printf("A")}；该 case 后没有 \code{break}，流程“穿透”到 \code{case 2}，再输出 B；\code{case 2} 后有 \code{break}，跳出 switch。因此输出 AB，不会执行 \code{case 3} 和 \code{default}。\end{analysis}

\q B。\begin{analysis}\code{break} 在 \code{switch} 中的作用是终止 switch 结构的执行，使流程转到 switch 语句之后；在循环中 \code{break} 则用于跳出本层循环。\end{analysis}

\q B。\begin{analysis}条件运算符先判断 \code{a > b}，即 3 $>$ 5 不成立，条件表达式的值为冒号后面表达式的值，即 \code{b} 的值 5。\end{analysis}

\q B。\begin{analysis}优先级由高到低为：\code{!}（单目逻辑非，最高）$>$ \code{>}（关系运算符）$>$ \code{\&\&}（逻辑与）$>$ \code{=}（赋值运算符），所以 \code{!} 的优先级最高。\end{analysis}

\q B。\begin{analysis}C 语言规定 \code{else} 总是与它前面最近的、尚未配对的 \code{if} 配对，与书写时的缩进无关。一个 \code{if} 最多只能有一个 \code{else}，\code{else} 也不能单独使用。\end{analysis}

\q A。\begin{analysis}闰年的判断条件是：能被 4 整除且不能被 100 整除，或者能被 400 整除，即 \code{(year \% 4 == 0 \&\& year \% 100 != 0) || year \% 400 == 0}。B 中 \code{year \% 400 = 0} 把 \code{==} 误写成 \code{=}；C 漏掉了被 400 整除的情况且把 \code{==} 写成 \code{=}；D 的逻辑关系错误。\end{analysis}

\q C。\begin{analysis}\code{switch} 后面圆括号内的表达式必须是整型（包括字符型、枚举型），不能是浮点型；\code{case} 后的常量表达式的值也必须是与该表达式类型一致的整型常量。\end{analysis}

\answerpart{二、多项选择题（5 题，每题 3 分，共 15 分）}

\q ABC。\begin{analysis}A、B 正确：关系表达式的值为 1 或 0，C 语言以非 0 为真、0 为假；C 正确：逻辑运算符优先级为 \code{!} $>$ \code{\&\&} $>$ \code{||}。D 错：\code{||} 的优先级低于 \code{\&\&}。\end{analysis}

\q ABC。\begin{analysis}A 对：\code{if} 后的条件可以是任意合法的表达式，按非 0 为真判断；B 对：条件必须用圆括号括起来；C 对：单分支 \code{if} 可以没有 \code{else}；D 错：\code{else} 后面跟的是语句（可以是复合语句），不能随意加分号，若写成 \code{else;} 则表示空语句。\end{analysis}

\q ABC。\begin{analysis}A、B、C 都是 switch 的语法规定：表达式必须为整型或字符型；case 后必须是常量表达式；各 case 常量值不能相同。D 错：\code{default} 分支不是必需的，可以省略。\end{analysis}

\q ABC。\begin{analysis}A 对：\code{?:} 是 C 语言唯一的三目运算符；B 对：条件运算符是右结合的；C 对：表达式1 的值为非 0 时取表达式2 的值。D 错：表达式 2 与表达式 3 的类型不要求完全相同，二者类型不同时会按自动类型转换规则处理，结果类型为二者转换后的共同类型。\end{analysis}

\q ABC。\begin{analysis}A 对：\code{x \% 2 == 0} 表示 x 除以 2 的余数为 0，即 x 是偶数；B 对：\code{!(x \% 2)} 中 x 为偶数时 \code{x \% 2} 为 0，逻辑非得 1；C 对：x 为偶数时 \code{x / 2 * 2} 等于 x。D 错：\code{x \% 2 = 0} 是把 0 赋给 \code{x \% 2}，语法错误（赋值号左边不是变量）。\end{analysis}

\answerpart{三、判断题（5 题，每题 2 分，共 10 分）}

\q √。\begin{analysis}C 语言没有专门的逻辑类型，\code{if} 的条件按“非 0 为真、0 为假”判断，因此条件可以是任意合法的表达式，如 \code{if (a)}、\code{if (a = 5)}。\end{analysis}

\q √。\begin{analysis}这是 C 语言关于悬空 else 的规定：\code{else} 与它前面最近的、尚未配对的 \code{if} 配对，书写缩进、换行都不影响配对关系，必要时应使用花括号改变配对。\end{analysis}

\q ×。改正：\code{switch} 语句中每个 \code{case} 后不一定需要 \code{break}；若不加 \code{break}，流程会继续执行后面 case 的语句（穿透）。\begin{analysis}\code{break} 用于跳出 switch；多个 case 共用同一段处理语句时，前面的 case 可以故意不加 break。是否加 break 取决于程序逻辑。\end{analysis}

\q ×。改正：\code{3 < x < 10} 的运算顺序是 \code{(3 < x) < 10}，前一个关系表达式的值为 1 或 0，不能表示 $3<x<10$；应写成 \code{x > 3 \&\& x < 10}。\begin{analysis}关系运算符是左结合的，\code{3 < x} 成立时得 1，再判断 \code{1 < 10} 恒为真，因此 \code{3 < x < 10} 永远为真（x 小于等于 3 时得 0，0 $<$ 10 也为真），与数学区间的含义完全不同。\end{analysis}

\q √。\begin{analysis}优先级由高到低为：算术运算符 $>$ 关系运算符 $>$ 逻辑运算符 $>$ 条件运算符 $>$ 赋值运算符，所以 \code{a + 1 > b} 会先算 \code{a + 1}，关系运算符又高于赋值运算符。\end{analysis}

\answerpart{四、填空题（10 空，每空 1 分，共 10 分）}

\q （1）1；（2）0。\begin{analysis}关系表达式成立时值为 1，不成立时值为 0；这是 C 语言中表示逻辑真、假的方式。\end{analysis}

\q （3）5。\begin{analysis}\code{a > b} 即 3 $>$ 5 不成立，条件运算符取冒号后表达式的值，即 \code{b} 的值 5。\end{analysis}

\q （4）假（0）；（5）短路。\begin{analysis}逻辑与 \code{\&\&} 的左边为假时，整个表达式已确定为假，右边的表达式不再计算，这种特性称为短路求值。\end{analysis}

\q （6）\code{break}。\begin{analysis}\code{break} 使流程跳出 switch 结构；若不加 break，会继续执行后面 case 的语句。\end{analysis}

\q （7）\code{?:}。\begin{analysis}条件运算符由问号和冒号组成，形式为 \code{表达式1 ? 表达式2 : 表达式3}，是 C 语言唯一的三目运算符。\end{analysis}

\q （8）\code{x \% 2 != 0}（或 \code{x \% 2 == 1}）。\begin{analysis}奇数除以 2 的余数为 1，因此条件可写为 \code{x \% 2 != 0}；也可以用 \code{x \% 2} 直接作为条件（非 0 即真）。\end{analysis}

\q （9）1；（10）1。\begin{analysis}\code{x} 为 0 时 \code{!x} 为逻辑非，结果为 1；\code{x == 0} 成立，关系表达式的值也为 1。\end{analysis}

\answerpart{五、程序阅读题（2 题，每题 5 分，共 10 分）}

\q 运行结果为：
\begin{lstlisting}
a=1,b=1,c=0
\end{lstlisting}
\begin{analysis}\code{c = a++ \&\& b++;} 先计算左边 \code{a++}：取 \code{a} 的当前值 0 作为左操作数，\code{a} 自增为 1。由于左操作数为 0（假），逻辑与发生短路，右边的 \code{b++} 不再计算，\code{b} 保持为 1。逻辑与的结果为 0，赋给 \code{c}，所以 \code{c = 0}。输出 \code{a=1,b=1,c=0}。\end{analysis}

\q 运行结果为：
\begin{lstlisting}
23
\end{lstlisting}
\begin{analysis}\code{k = 2}，与 \code{case 2} 匹配，执行 \code{printf("2")}；\code{case 2} 后没有 \code{break}，流程穿透到 \code{case 3}，紧接输出 3，然后遇到 \code{break} 跳出 switch；\code{default} 不执行。因此输出 23（case 1 被跳过）。\end{analysis}

\answerpart{六、简答题（3 题，每题 5 分，共 15 分）}

\q 参考答案要点：
（1）优先级由高到低：逻辑非 \code{!}（单目）$>$ 算术运算符 $>$ 关系运算符（其中 \code{<}、\code{<=}、\code{>}、\code{>=} 高于 \code{==}、\code{!=}）$>$ 逻辑与 \code{\&\&} $>$ 逻辑或 \code{||} $>$ 条件运算符 $>$ 赋值运算符 $>$ 逗号运算符（2 分）。
（2）关系运算符和逻辑运算符的运算结果都是 1（真）或 0（假），运算结果可以直接作为 \code{if}、\code{while} 等的条件（1 分）。
（3）短路求值：\code{A \&\& B} 中，若 A 为假，则不再计算 B，整个表达式为假；\code{A || B} 中，若 A 为真，则不再计算 B，整个表达式为真（1 分）。
（4）利用短路可以简化条件、提高效率并避免错误，例如 \code{if (a != 0 \&\& b / a > 1)} 中，当 \code{a} 为 0 时不会执行 \code{b / a}，从而避免除数为 0（1 分）。

\q 参考答案要点：
（1）单分支 \code{if} 语句：\code{if (条件) 语句}。条件成立时执行后面的语句，否则不执行（1 分）。
（2）双分支 \code{if-else} 语句：\code{if (条件) 语句1 else 语句2}。条件成立执行语句 1，否则执行语句 2，二者必执行其一（2 分）。
（3）多分支 \code{if……else if……else} 结构：\code{if (条件1) 语句1 else if (条件2) 语句2 …… else 语句n}。从上到下依次判断，执行第一个成立的条件对应的语句，其余分支都不执行；若都不成立则执行最后的 \code{else}（没有 \code{else} 则什么都不执行）（2 分）。

\q 参考答案要点：
（1）执行过程：先计算 switch 后括号内表达式的值，然后依次与各 case 后的常量比较，找到相等的 case 便从该处开始执行；若某 case 后有 \code{break}，执行到 break 即跳出 switch；若没有 break 则继续执行后续 case 的语句（穿透）；若没有匹配的 case 且有 \code{default}，则执行 default 后的语句（2 分）。
（2）区别：\code{if} 语句可判断任意表达式（区间、逻辑组合等），\code{switch} 只能对整型或字符型表达式做等值判断（2 分）。
（3）适用场合：分支较少或需要判断区间、复杂条件时用 \code{if}；分支较多且都是对某个整型或字符型表达式的等值判断时，用 \code{switch} 更清晰（1 分）。

\answerpart{七、综合题（1 题，共 20 分）}

\q 参考程序如下：
\begin{lstlisting}[language={[custom]C}]
#include <stdio.h>
int main(void)
{
    int year;
    scanf("%d", &year);
    if ((year % 4 == 0 && year % 100 != 0) || year % 400 == 0)
        printf("%d is a leap year\n", year);
    else
        printf("%d is not a leap year\n", year);
    return 0;
}
\end{lstlisting}
\begin{analysis}运行结果（输入 2000、1900、2024 时）为：
\begin{lstlisting}
2000 is a leap year
1900 is not a leap year
2024 is a leap year
\end{lstlisting}
评分要点：正确包含头文件并定义整型变量 \code{year}（3 分）；\code{scanf} 正确读入年份（3 分）；闰年条件完整正确（能被 4 整除且不能被 100 整除，或能被 400 整除），括号位置正确（8 分）；\code{if-else} 两个分支的输出正确（4 分）；程序完整可编译、有 \code{return 0;}（2 分）。\end{analysis}
\begin{analysis}易错提示：闰年条件中“能被 4 整除且不能被 100 整除”必须用括号括起来再与“能被 400 整除”相“或”，否则 \code{\&\&} 优先级高于 \code{||} 会使逻辑关系出错；若漏掉“能被 400 整除”则 2000 会被误判为非闰年。\end{analysis}

\end{document}
